
\documentclass[12pt, a4paper]{article}

\usepackage[T1]{fontenc}
\usepackage[utf8]{inputenc}
\usepackage[polish]{babel}
\usepackage{lmodern}
\usepackage{microtype}

\usepackage{amsmath, amssymb, amsthm}
\usepackage{mathtools}

\usepackage{indentfirst}
\setlength{\parindent}{0.8cm}

\usepackage{graphicx}
\usepackage{booktabs}
\usepackage{enumitem}
\usepackage{url}

\usepackage{geometry}
\geometry{
    a4paper,
    margin=2.5cm
}

\usepackage[
    colorlinks=true,
    linkcolor=black,
    citecolor=black,
    urlcolor=blue
]{hyperref}

\theoremstyle{definition}
\newtheorem{definicja}{Definicja}[section]

\theoremstyle{plain}
\newtheorem{twierdzenie}{Twierdzenie}[section]
\newtheorem{lemat}[twierdzenie]{Lemat}

\theoremstyle{remark}
\newtheorem{uwaga}{Uwaga}[section]

\title{\textbf {\Huge{Projekt inżynierski}}\\ \vspace{3pt}
    \Large{{Rozpoznawanie instrumentów muzycznych
    \\ w utworach polifonicznych z wykorzystaniem sieci neuronowych}}}

\author{
    Kacper Haczkiewicz,
    Wiktor Niedźwiedzki,
    Dawid Skowroński
}

\date{}

\begin{document}

\maketitle

\tableofcontents
\newpage

\section{Wprowadzenie}

Niniejszy projekt inżynierski poświęcony jest problemowi klasyfikacji instrumentów muzycznych w utworach polifonicznych. W tego rodzaju utworach dźwięki pochodzące z różnych instrumentów mogą występować jednocześnie, co utrudnia ich prawidłową identyfikację. Do rozwiązania tego problemu można wykorzystać metody analizy sygnałów oraz metody uczenia maszynowego oparte na sieciach neuronowych.

W projekcie wykorzystamy krótkoczasową transformatę Fouriera (STFT) do analizy sygnału audio i wyznaczenia spektrogramu melowego, który posłuży jako dane wejściowe dla konwolucyjnej sieci neuronowej. Na podstawie uzyskanych wyników ocenimy skuteczność zastosowanego rozwiązania.


\section{Cel i zakres projektu}

\subsection{Cel projektu}

Celem projektu jest implementacja i ocena skuteczności konwolucyjnej sieci neuronowej (CNN) w rozpoznawaniu instrumentów muzycznych w utworach polifonicznych. W tym celu sygnał audio zostanie przekształcony z dziedziny czasu do reprezentacji czasowo-częstotliwościowej w postaci mel-spektrogramu z wykorzystaniem krótkoczasowej transformaty Fouriera (STFT). Tak przygotowane dane zostaną wykorzystane do trenowania i testowania modelu.

\subsection{Zakres projektu}

Tu napiszemy, jak bardziej będzie wiadomo co i jak (chyba że wiecie jak to lepiej napisać to napiszćie)

\newpage

\section{Podstawy teoretyczne}

\subsection{Sygnał audio i jego reprezentacja cyfrowa}


Jak podaje Zieliński \cite{zielinski2005}, sygnał można rozumieć jako zmienność wielkości fizycznej, którą opisuje się za pomocą funkcji. W przypadku dźwięku jest to zmiana ciśnienia akustycznego w czasie. Podczas nagrywania mikrofon przekształca falę akustyczną w sygnał elektryczny, który może zostać zapisany i poddany dalszej analizie.

Komputer przetwarza dźwięk w postaci cyfrowej, reprezentowanej przez ciąg próbek, czyli wartości sygnału pobrane w określonych chwilach czasu. Proces ten obejmuje próbkowanie oraz kwantyzację. Próbkowanie polega na pobieraniu kolejnych wartości sygnału, natomiast kwantyzacja przyporządkowuje im określone poziomy amplitudy \cite{zielinski2005,oppenheim2010}.

Liczbę próbek pobieranych w ciągu sekundy określa częstotliwość próbkowania $f_s$, wyrażana w hercach (Hz). Jest ona odwrotnością okresu próbkowania $T_s$:

\begin{equation}
    f_s = \frac{1}{T_s}.
\end{equation}

Odpowiedni dobór częstotliwości próbkowania pozwala zachować informacje o składowych częstotliwościowych sygnału. Zbyt niska częstotliwość może prowadzić do aliasingu, czyli błędnego odwzorowania składowych widma. Aby ograniczyć to zjawisko, w procesie konwersji analogowo-cyfrowej stosuje się filtr antyaliasingowy \cite{zielinski2005, oppenheim2010}.

\subsection{Analiza częstotliwościowa sygnałów}

Analiza częstotliwościowa umożliwia opisanie sygnału poprzez jego składowe częstotliwościowe. Jej podstawą jest rozkład sygnału na funkcje harmoniczne (\textit{funkcje sinusoidalne}), który pozwala określić, jakie częstotliwości występują w sygnale oraz jaki jest ich udział \cite{zielinski2005,oppenheim2010}.

\textbf{Sygnał} można traktować jako element odpowiedniej przestrzeni matematycznej. W przypadku cyfrowego nagrania skończony fragment sygnału jest wektorem próbek $x[n]$. Jego normę euklidesową definiuje się jako:

\begin{equation}
    \|x\|_2 =
    \sqrt{\sum_{n=0}^{N-1}|x[n]|^2},
\end{equation}

gdzie $N$ oznacza liczbę próbek. Norma pozwala określić wielkość sygnału, natomiast iloczyn skalarny umożliwia opisywanie zależności pomiędzy sygnałami. Pojęcia te są wykorzystywane w analizie funkcji ortogonalnych i rozwinięć Fouriera. Dla nieskończonych ciągów o sumowalnym kwadracie modułu analogiczną konstrukcję stanowi przestrzeń $\ell^2$ \cite{zielinski2005}.

\textbf{Dyskretna transformata Fouriera} (DFT) przekształca skończony ciąg próbek w zestaw współczynników opisujących jego składowe częstotliwościowe. Definiuje się ją wzorem:

\begin{equation}
    X[k] = \sum_{n=0}^{N-1}
    x[n]e^{-j\frac{2\pi}{N}kn},
    \quad k=0,\ldots,N-1,
\end{equation}

gdzie $x[n]$ oznacza próbki sygnału, a $X[k]$ jest zespolonym współczynnikiem odpowiadającym indeksowi częstotliwościowemu $k$. Moduł współczynnika opisuje poziom danej składowej, natomiast jego argument określa fazę. Częstotliwość odpowiadającą indeksowi $k$ wyznacza zależność:

\begin{equation}
    f_k = \frac{k f_s}{N},
\end{equation}

gdzie $f_s$ oznacza częstotliwość próbkowania. Szybka transformata Fouriera (FFT) jest rodziną algorytmów umożliwiających efektywne obliczanie DFT \cite{oppenheim2010}.

\textbf{Widmo amplitudowe} przedstawia moduły współczynników DFT, natomiast widmo mocy jest związane z kwadratami ich modułów. Obie reprezentacje pozwalają analizować rozkład składowych częstotliwościowych sygnału. Ponieważ dźwięk może zmieniać się w czasie, analizę często przeprowadza się dla kolejnych fragmentów sygnału. Służy do tego krótkoczasowa transformata Fouriera (STFT), w której każdy fragment jest analizowany osobno, zwykle po przemnożeniu go przez funkcję okna ograniczającą efekty wynikające z jego odcięcia.

Wyniki STFT można przedstawić w postaci \textbf{spektrogramu}, czyli reprezentacji pokazującej zmiany zawartości częstotliwościowej sygnału w czasie.

\subsection{Skala melowa i mel-spektrogram}

Skala melowa jest nieliniową skalą częstotliwości, zaprojektowaną w celu przybliżenia sposobu, w jaki człowiek postrzega wysokość dźwięku. W skali tej częstotliwości niższe są rozróżniane z większą dokładnością niż częstotliwości wyższe. Przeliczenie częstotliwości wyrażonej w hercach na skalę melową opisuje zależność \cite{fayek2016}:

\begin{equation}
    m(f) = 2595 \log_{10}\left(1+\frac{f}{700}\right),
\end{equation}
gdzie $f$ oznacza częstotliwość w hercach, a $m(f)$ jej wartość w melach.

Mel-spektrogram otrzymuje się przez zastosowanie banku filtrów melowych do widma mocy kolejnych ramek sygnału. Bank składa się z filtrów trójkątnych rozmieszczonych równomiernie w skali melowej. Każdy filtr przypisuje prążkom widma wagi zależne od ich położenia, a następnie agreguje ich wartości w obrębie danego pasma \cite{fayek2016}.

Wartość $m$-tego pasma dla ramki $t$ można zapisać jako:

\begin{equation}
    S_m[t] = \sum_{k=0}^{K-1} H_m[k]P_t[k],
\end{equation}

gdzie $P_t[k]$ oznacza wartość widma mocy dla $k$-tego prążka w ramce $t$, natomiast $H_m[k]$ jest wagą $m$-tego filtru melowego. Otrzymana macierz $S_m[t]$ przedstawia rozkład energii w pasmach melowych w kolejnych chwilach czasu. Jej wizualizacją jest mel-spektrogram, w którym osie odpowiadają czasowi i częstotliwości w skali melowej, a kolor reprezentuje wartość sygnału w danym paśmie.

\newpage
\section*{Bibliografia}
\addcontentsline{toc}{section}{Bibliografia}

\begin{thebibliography}{9}

\bibitem{zielinski2005}
T.~P. Zieliński,
\textit{Cyfrowe przetwarzanie sygnałów: od teorii do zastosowań}.
Warszawa: Wydawnictwa Komunikacji i Łączności, 2005.

\bibitem{oppenheim2010}
A.~V. Oppenheim i R.~W. Schafer,
\textit{Discrete-Time Signal Processing},
wyd. 3. Pearson, 2010.

\bibitem{fayek2016}
H. M. Fayek, \textit{Speech Processing for Machine Learning: Filter banks, Mel-Frequency Cepstral Coefficients (MFCCs) and What's In-Between}, 2016.
Dostęp: \url{https://haythamfayek.com/2016/04/21/speech-processing-for-machine-learning.html}

\end{thebibliography}

\end{document}
